\documentclass[11pt]{article}
\usepackage[a4paper,margin=27mm]{geometry}
\usepackage{amsmath,amssymb,amsthm,hyperref}
\newcommand{\MA}{\operatorname{MA}}
\newcommand{\NA}{\mathrm{NA}}
\newcommand{\ord}{\operatorname{ord}}
\newtheorem{theorem}{Theorem}
\newtheorem{lemma}[theorem]{Lemma}
\title{Normalized model recovery can lose a singularity endpoint}
\author{}
\date{6 October 2026}
\begin{document}
\maketitle
\begin{abstract}
We construct a candidate counterexample to endpoint continuity under
normalized model energy and entropy recovery. All configurations are
actual smooth equivariant blowups of $\mathbf P^1\times\mathbf P^1$,
with relatively ample rational bundles. Their data converge uniformly
to zero and their non-Archimedean energies and entropies converge to
zero. Every ray starts at zero and every cutoff remains zero. Nevertheless
the endpoint $P$-model potential has raw mass tending to zero, whereas
the target endpoint has mass one. Every fixed strictly interior level
is eventually recovered exactly. This diagnoses a missing implication;
it does not disprove existence of a better approximation for P19.2.
The construction has not yet passed independent mathematical review.
\end{abstract}

\section{Conventions and statement}
Use $X=\mathbf P^1_{\mathbf C}$,
$\omega=dd^c\log(1+|z|^2)$ of mass one, and
\[
 dd^c=\frac{i}{2\pi}\partial\bar\partial,\qquad
 t=-\log|\tau|^2 .                                      \tag{1}
\]
The analytic measure is the non-pluripolar curvature, probability
normalized only for finite-energy ray slices, and $H$ is full probability
entropy against $\omega$. Level masses are raw, including zero.
For an actual rational vertical divisor $D$ on a dominating configuration,
its usual datum at a Gauss component is $f_D(v_F)=D_F/b_F$,
where $b_F=\ord_F(\tau)$. On $X^\beth$ over trivially valued
$\mathbf C$ retain all semivaluation supports.

\begin{theorem}[Candidate diagnostic]
There exist smooth equivariant projective configurations $\mathcal Y_j$
with relatively ample rational bundles in class $\mathcal O(1)$ and
actual zero-start weak rays $u_j$, with data $\phi_j$, such that
\[
 \sup\phi_j=0,\quad \|\phi_j\|_\infty\to0,\quad
 E_\NA(\phi_j)\to0,\quad H_\NA(\phi_j)\to0,
 \quad H(u_{j,t})/t\to0\quad(j,t\to\infty).               \tag{2}
\]
Their cutoffs are all zero. Every fixed inverse Legendre level
$\lambda<0$ is eventually exactly zero. At $\lambda=0$ their genuine
$P$-model levels have raw mass $1/j$, while the target zero ray has
endpoint zero and mass one. Endpoint poles are not recovered.
\end{theorem}

\section{Actual configurations and positivity}
Fix an integer $m\geq2$. Beginning with
$\mathbf P^1_z\times\mathbf P^1_\tau$, blow up $(z,\tau)=(0,0)$.
Inductively blow up the intersection of the newest exceptional component
and the strict transform of $\tau=0$, until there are $m$ exceptional
components $F_1,\ldots,F_m$. Write $F_0$ for the strict central component.
These centers are smooth fixed points of the $\mathbf C^*$ action on
$\tau$, so the resulting compact surface $\mathcal Y_m$ is smooth,
projective and equivariant. In the $(\ord(z),\ord(\tau))$ lattice the
exceptional rays are $(1,i)$, $1\leq i\leq m$. Adjacent determinants
are one. Thus
\[
 b_0=1,\quad b_i=i,\quad v_{F_i}=i^{-1}\ord_0,\quad
 B=F_0+\sum_{i=1}^m iF_i .                              \tag{3}
\]
The central incidence chain is $F_0,F_m,F_{m-1},\ldots,F_1$.
Here $F_0^2=-m$, $F_i^2=-2$ for $i<m$, and $F_m^2=-1$;
adjacent central intersections are one, others zero. These assertions
follow inductively from the blowup self-intersection formula, retaining
the multiplicities in (3).

Choose rational $a_i>0$ with $c=\sum_i a_i<1$, and put
\[
 d_i=-\sum_{k=1}^m a_k\min(1,i/k),\qquad
 D=\sum_{i=1}^m d_iF_i,\qquad
 \mathcal L=p_z^*\mathcal O(1)+D.                        \tag{4}
\]
This is an actual rational line bundle: clear its finitely many
denominators. Direct intersection on the chain gives
\[
 \mathcal L.F_0=1-c>0,\qquad
 \mathcal L.F_i=a_i/i>0\quad(1\leq i\leq m).             \tag{5}
\]
For $i<m$ use $D.F_i=-2d_i+d_{i-1}+d_{i+1}$ with $d_0=0$;
the summand $-a_k\min(1,i/k)$ has second difference $a_k/k$ exactly
at $i=k$. For $i=m$, $D.F_m=-d_m+d_{m-1}=a_m/m$.
For $F_0$ use $D.F_0=d_m=-c$ and $p_z^*\mathcal O(1).F_0=1$.

All other fibers have degree one. A line bundle on a proper curve is
ample exactly when it has positive degree on every irreducible component;
ampleness is unchanged upon reduction of a nilpotent thickening. Apply
this to the nonreduced fiber (3), after clearing denominators.
The proper-morphism fiber criterion [Stacks, Lemma 37.50.3] therefore
makes $\mathcal L$ relatively ample. Twisting by a sufficiently positive
base bundle gives an ample bundle on the compact surface and hence
a K\"ahler representative. Consequently this is an admissible smooth
relatively K\"ahler cohomological configuration, including its actual
nonreduced central fiber. No formal fractional-speed configuration is
substituted for these blowups.

\section{The datum and its actual weak ray}
Write $v_i=1/i$, $S_0=0$, $S_i=\sum_{k=1}^i a_k$, and
$r=\sum_i a_iv_i$. Define a rational convex piecewise affine function
$g:[0,1]\to\mathbf R_{\geq0}$ by
\[
 g(x)=\int_x^c v(s)\,ds\quad(0\leq x\leq c),\qquad
 g(x)=0\quad(c\leq x\leq1),\qquad
 v(s)=v_i\text{ on }(S_{i-1},S_i).                       \tag{6}
\]
Its slopes are $-v_1,\ldots,-v_m,0$, in increasing order,
and $g(0)=r$. Elementary one-variable convex duality gives, for $v\geq0$,
\[
 \sup_{0\leq x\leq1}\{-xv-g(x)\}
             =-\sum_i a_i\min(v,v_i).                   \tag{7}
\]
At $v=1/k$ the right side is $d_k/k$, exactly (4). It is zero
on the trivial branch and on branches centered at other points,
including infinity in its local $\mathcal O(1)$ frame. At the supported
semivaluation $z=0$ it equals $-r$, without an undefined difference.
Thus the actual model datum is
$\phi(v)=-\sum_i a_i\min(v,v_i)$ on the zero branch and zero elsewhere;
in particular $\sup\phi=0$ and $\|\phi\|_\infty=r$.

Let $F(x)=x\log x+(1-x)\log(1-x)$ with $0\log0=0$.
For $\rho=\log|z|^2$ set
\[
 q_t(\rho)=\sup_{0\leq x\leq1}\{x\rho-F(x)-tg(x)\},
 \qquad u_t=q_t(\rho)-\log(1+e^\rho).                  \tag{8}
\]
Then $u_0=0$ and $-rt\leq u_t\leq0$. The functions extend across zero
and infinity at every finite time. Their total weight is compatible
with the actual divisor (4). To check bounded compatibility, omit $F$
in (8) and call the resulting support function $h$; since
$-\log2\leq F\leq0$, $h\leq q\leq h+\log2$ uniformly.
The maximum defining $h$ is attained at the finitely many nodes $S_i$
and at $1$. With a common denominator $M$ it is the logarithmic maximum
of the actual monomials
\[
 z^{MS_i}\tau^{Mg(S_i)}\ (0\leq i\leq m),\qquad z^M.     \tag{9}
\]
Their minimum order on $F_k$ is
$M\min_x(x+kg(x))=-Md_k$. On the smooth fan of the preceding blowups
their support function is linear on each cone: its break directions
are precisely $(1,i)$. Hence their common ideal becomes invertible
on this model. Removing its canonical factor removes exactly $-MD$.
The residual logarithmic maximum is locally bounded in each actual
bundle frame. The uniform comparison with $q$ proves
$U+\log|s_D|^2=O(1)$, with the same logarithm-square convention.

For completeness we identify the actual weak ray without assuming that
an arbitrary subgeodesic is geodesic. Its inverse Legendre curve is
\[
 \chi(\lambda)=\inf_{t\geq0}(u_t-t\lambda)=
 \begin{cases}
 0,&\lambda\leq-r,\\
 w_{b(\lambda)},&-r<\lambda<0,\quad g(b(\lambda))=-\lambda,\\
 w_c,&\lambda=0,\\
 -\infty,&\lambda>0,
 \end{cases}                                           \tag{10}
\]
For the interior formula maximize $x\rho-F(x)$ under $g(x)\leq-\lambda$.
Convex Lagrange duality applies since $x=1$ strictly satisfies the
constraint for $\lambda<0$; equivalently use elementary convex duality
on the finite affine pieces. The feasible interval is $[b(\lambda),1]$.
At zero let $t\to\infty$ in (8). Compactness of maximizing $x$ forces
all cluster points into $[c,1]$, proving the endpoint formula as well.
The functions in (10) are
\[
 w_b(z)=
 \begin{cases}
 b\log|z|^2-F(b)-\log(1+|z|^2),&|z|^2\leq b/(1-b),\\
 0,&|z|^2\geq b/(1-b),
 \end{cases}\qquad 0<b<1 .                             \tag{11}
\]
They are genuine normalized $P$-model potentials. Indeed for any rooftop
competitor $h\leq0$, $h\leq w_b+C$, its local weight
$Q=h+\log(1+|z|^2)$ satisfies that $Q-b\log|z|^2$ is subharmonic
and bounded above near zero, hence extends there. On the circle
$|z|^2=b/(1-b)$, $Q\leq\log(1+b/(1-b))$.
The maximum principle gives $h\leq w_b$ on the disk; outside it
$h\leq0=w_b$. Since $w_b$ itself is a competitor, $P_\omega[w_b]=w_b$.
Its curvature has the polar atom $b\delta_0$ and agrees with $\omega$
outside that disk, with no boundary-circle atom because the first
derivatives match. Its raw non-pluripolar mass is therefore $1-b$.

Curve (10) is parameter-concave and upper semicontinuous, being the
Legendre infimum in (8), constant zero below $-r$ and $P$-maximal at
every nontrivial level, including zero. It is a bounded maximal test
curve. The analytic maximal-test-curve/weak-ray correspondence [DXZ, DZ]
therefore makes its inverse Legendre transform the actual bounded
weak ray from zero. Elementary convex biconjugacy identifies this
transform with (8). All finite time slices are in $\mathcal E^1$.
This uses the analytic K\"ahler correspondence in its initial-zero
signature and no polarization-dependent closure theorem.

\section{Energy, entropy, and endpoint separation}
The raw model probability measure, retaining the multiplicities (3), is
\[
 \MA_\NA(\phi)=(1-c)\delta_{v_{\rm triv}}
                       +\sum_{i=1}^m a_i\delta_{v_i\ord_0}.
 \tag{12}
\]
As $A_X(v_i\ord_0)=v_i$, its discrepancy entropy is $r$.
The intersection energy and actual ray energy are
\[
 E_\NA(\phi)=D^2/2
    =-\frac12\sum_{i,k}a_ia_k\min(v_i,v_k)
    =-\int_0^1g(x)\,dx,
 \qquad E(u_t)=-t\int_0^1g(x)\,dx .                    \tag{13}
\]
The first equality follows by intersecting $D$ with (5);
the elementary integral in (6) gives the double sum.
For the last identity the radial probability is $dx$, where
$x=\partial_\rho q_t$, and $\partial_tu_t=-g(x)$ at its maximizing
point. The bounded-geodesic energy derivative thus gives (13), starting
from $E(0)=0$. Values at the finitely many piece boundaries have no
radial probability mass. In particular $|E_\NA(\phi)|\leq r$.

On $x\in(S_{i-1},S_i)$, the maximizing relation is
$\rho=\log(x/(1-x))-tv_i$. The density against $\omega$ is
\[
 f_t=e^{tv_i}(1-x+e^{-tv_i}x)^2.                         \tag{14}
\]
The plateau intervals of $q_t$ in $\rho$ have zero curvature mass;
on $x\in(c,1)$ the density is one. First derivatives match at all
interfaces, so there is no missing circle mass. It follows directly
that every finite-time entropy is finite and
\[
 H(u_t)=tr+2\sum_i\int_{S_{i-1}}^{S_i}
                       \log(1-x+e^{-tv_i}x)\,dx,
 \qquad -2\leq H(u_t)-tr\leq0.                          \tag{15}
\]
The last bound uses $\log(1-x)\leq\log(1-x+e^{-tv_i}x)\leq0$
and $\int_0^1-\log(1-x)dx=1$. It is uniform in the changing model.
Consequently the analytic entropy slope is exactly $r$, matching (12).

Now take $m=j\geq2$ and
\[
 a_i=j^{-3}\ (i<j),\qquad
 a_j=1-j^{-1}-(j-1)j^{-3},\qquad c_j=1-j^{-1}.           \tag{16}
\]
All coefficients in (16) are positive. Also
\[
 r_j=\frac{a_j}{j}+j^{-3}\sum_{i=1}^{j-1}\frac1i\to0.
\]
Equations (7), (12)--(15) prove (2); uniform model convergence implies
strong convergence, also directly consistent with (13).
For any fixed $\lambda<0$, eventually $r_j<-\lambda$, and (10) gives
$\chi_j(\lambda)=0$ exactly. The cutoff remains zero at every index.
Nevertheless $\chi_j(0)=w_{c_j}$ has raw mass $1-c_j=1/j$ and
a pole at zero of coefficient $c_j$, while the zero target ray has
$\chi(0)=0$ of mass one and no pole. In dimension one the ordered
singularity metric is $d_{\mathcal S}([w_{c_j}],[0])=c_j/2$
with probability-normalized energy [DDNL, Lemma 3.4]; it stays away
from zero. This is a failure at the normalized, unchanged cutoff itself,
rather than a constant translation of levels.

\section{The original problem remains an existence question}
These configurations refute the implication that normalized strong
energy and discrepancy-entropy recovery, even with a uniform joint
analytic entropy slope estimate and eventual exact recovery of every
fixed interior level, must recover the endpoint $P$ data. They satisfy
actual relative ampleness rather than only a limiting nef class.
The target datum zero also has a constant trivial configuration
sequence preserving its endpoint exactly. Thus this construction cannot
be a counterexample to existence of the simultaneous approximation
asked for in P19.2. It establishes a specific diagnostic obstruction
and leaves the full residual target open.

\begin{thebibliography}{9}
\bibitem{Stacks}[Stacks] The Stacks Project,
\href{https://stacks.math.columbia.edu/tag/0D2S}{Lemma 37.50.3,
Tag 0D2S}, relative ampleness from an ample proper fiber; the usual
component-degree criterion for proper curves and nilpotent invariance
of ampleness are applied to the actual fiber (3).
\bibitem{DXZ}[DXZ] T. Darvas, M. Xia and K. Zhang,
\emph{A transcendental approach to non-Archimedean metrics of
pseudoeffective classes}, Comment. Math. Helv. 100 (2025), 269--322,
Theorem 2.6 in its analytic initial-zero K\"ahler signature.
\bibitem{DZ}[DZ] T. Darvas and K. Zhang,
\href{https://arxiv.org/pdf/2208.08324v4}{arXiv:2208.08324v4},
Theorems 3.7 and 3.9 and the bounded analytic Legendre proofs.
\bibitem{DDNL}[DDNL] T. Darvas, E. Di Nezza and C. H. Lu,
\emph{The metric geometry of singularity types},
\href{https://arxiv.org/pdf/1909.00839}{primary text},
Lemma 3.4, with probability energy and the operator (1).
\end{thebibliography}
\end{document}
